\section{Statistical Methods}
\label{sec:stats}

Outcomes are paired binary, so we use McNemar's test \cite{mcnemar1947},
which Dietterich \cite{dietterich1998} identifies as the appropriate
choice for algorithms evaluated once per unit. We use the \textbf{exact
two-sided binomial form} throughout rather than the $\chi^2$
approximation, because several comparisons have few discordant pairs ---
the C-vs-baseline contrast has only $b+c=7$. Discordant counts $(b,c)$
accompany every $p$-value, since the test is a statement about those
counts alone. Effect sizes are absolute percentage-point differences with
95\% percentile bootstrap intervals, $B=2000$ resamples at fixed seed 0,
resampling the \emph{per-pair differences} so the pairing is preserved
\cite{efron1993}.

We pre-declare \textbf{Full vs.\ Baseline} as the primary comparison. The
confirmatory family comprises the four arm-vs-baseline tests on the
canonical grid, the planned B-vs-C contrast, and the four
adaptive-arm-vs-matched-control tests --- a family of nine, to which
Holm's sequentially rejective procedure \cite{holm1979} is applied
(\SItables{}). The held-out instances are analysed identically
but are \emph{not} pooled into the family: they are replications of an
already pre-declared comparison rather than new hypotheses, so their
$p$-values are reported uncorrected. Full procedural detail is in
\SItables{}.

One consequence deserves flagging here rather than in a footnote:
\textbf{module C's improvement over baseline does not survive the
correction} at this family size ($p=1.6\times10^{-2}$ against a threshold
of $1.25\times10^{-2}$). A smaller five-comparison family covering the
canonical ablation alone would have retained it. We report the larger
family because it is the one the completed validation actually tests.

Following Dietterich, McNemar's test speaks to the difference in error
proportion on \emph{one} evaluation set; two repeats per grid cell do not
characterise variability across independent re-runs of the whole
campaign, and no result here should be read as doing so.
